% !TeX program = pdflatex
\documentclass[11pt,a4paper]{amsart}

\newcommand{\notelanguage}{finnish}
% Shared preamble for course notes and exercise sets.

\usepackage[T1]{fontenc}
\usepackage{lmodern}
\providecommand{\notelanguage}{english}
\usepackage[finnish,english]{babel}
\AtBeginDocument{\expandafter\selectlanguage\expandafter{\notelanguage}}

\newcommand{\notesfinnishlanguage}{finnish}
\ifx\notelanguage\notesfinnishlanguage
\newcommand{\notestheoremname}{Lause}
\newcommand{\noteslemmaname}{Lemma}
\newcommand{\notespropositionname}{Propositio}
\newcommand{\notescorollaryname}{Korollaari}
\newcommand{\notesdefinitionname}{Määritelmä}
\newcommand{\notesexamplename}{Esimerkki}
\newcommand{\notesproblemname}{Tehtävä}
\newcommand{\notesexercisename}{Tehtävä}
\newcommand{\notesremarkname}{Huomautus}
\newcommand{\notessolutionname}{Ratkaisu}
\else
\newcommand{\notestheoremname}{Theorem}
\newcommand{\noteslemmaname}{Lemma}
\newcommand{\notespropositionname}{Proposition}
\newcommand{\notescorollaryname}{Corollary}
\newcommand{\notesdefinitionname}{Definition}
\newcommand{\notesexamplename}{Example}
\newcommand{\notesproblemname}{Problem}
\newcommand{\notesexercisename}{Exercise}
\newcommand{\notesremarkname}{Remark}
\newcommand{\notessolutionname}{Solution}
\fi

\usepackage[a4paper,margin=30mm]{geometry}
\usepackage{microtype}
\usepackage{graphicx}
\usepackage{enumitem}

\newcommand{\exerciseheading}[2]{%
  \par\noindent
  {\large\bfseries #1\par}
  \vspace{0.1em}
  \noindent{\large #2\par}
  \vspace{1.5em}
}

\usepackage{amsmath}
\usepackage{amssymb}
\usepackage{amsthm}
\usepackage{mathtools}
\usepackage{aliascnt}

% Row operations beside the affected row, e.g. \rowop{+2R_1}.
\newcommand{\rowop}[1]{\mathrel{\big|}\,#1}

% Linear systems: the argument gives the number of alignment pairs.
\newenvironment{linearsystem}[1]
{\left\{
  \begin{alignedat}{#1}}
    {
  \end{alignedat}\right.}

  \usepackage[hidelinks,pdfusetitle]{hyperref}

  \numberwithin{equation}{section}
  \setlist{itemsep=0.25em,topsep=0.5em}

  \theoremstyle{plain}
  \newtheorem{theorem}{\notestheoremname}[section]
  \newaliascnt{lemma}{theorem}
  \newtheorem{lemma}[lemma]{\noteslemmaname}
  \aliascntresetthe{lemma}
  \newaliascnt{proposition}{theorem}
  \newtheorem{proposition}[proposition]{\notespropositionname}
  \aliascntresetthe{proposition}
  \newaliascnt{corollary}{theorem}
  \newtheorem{corollary}[corollary]{\notescorollaryname}
  \aliascntresetthe{corollary}

  \theoremstyle{definition}
  \newaliascnt{definition}{theorem}
  \newtheorem{definition}[definition]{\notesdefinitionname}
  \aliascntresetthe{definition}
  \newaliascnt{example}{theorem}
  \newtheorem{example}[example]{\notesexamplename}
  \aliascntresetthe{example}
  \newaliascnt{problem}{theorem}
  \newtheorem{problem}[problem]{\notesproblemname}
  \aliascntresetthe{problem}
  \newtheorem{exercise}{\notesexercisename}

  \theoremstyle{remark}
  \newaliascnt{remark}{theorem}
  \newtheorem{remark}[remark]{\notesremarkname}
  \aliascntresetthe{remark}

  \newenvironment{solution}
  {
  \begin{proof}[\notessolutionname]}
    {
  \end{proof}}

  \usepackage[nameinlink,noabbrev]{cleveref}

  \crefname{theorem}{theorem}{theorems}
  \crefname{lemma}{lemma}{lemmas}
  \crefname{proposition}{proposition}{propositions}
  \crefname{corollary}{corollary}{corollaries}
  \crefname{definition}{definition}{definitions}
  \crefname{example}{example}{examples}
  \crefname{problem}{problem}{problems}
  \crefname{exercise}{exercise}{exercises}
  \crefname{remark}{remark}{remarks}

  \ifx\notelanguage\notesfinnishlanguage
  \crefname{theorem}{lause}{lauseet}
  \crefname{lemma}{lemma}{lemmat}
  \crefname{proposition}{propositio}{propositiot}
  \crefname{corollary}{korollaari}{korollaarit}
  \crefname{definition}{määritelmä}{määritelmät}
  \crefname{example}{esimerkki}{esimerkit}
  \crefname{problem}{tehtävä}{tehtävät}
  \crefname{exercise}{tehtävä}{tehtävät}
  \crefname{remark}{huomautus}{huomautukset}
  \fi

  \newcommand{\NN}{\mathbb{N}}
  \newcommand{\NNone}{\mathbb{N}_1}
  \newcommand{\ZZ}{\mathbb{Z}}
  \newcommand{\QQ}{\mathbb{Q}}
  \newcommand{\RR}{\mathbb{R}}
  \newcommand{\CC}{\mathbb{C}}
  \newcommand{\PP}{\mathbb{P}}

  \DeclareMathOperator{\im}{im}
  \DeclareMathOperator{\rank}{rank}
  \DeclareMathOperator{\syt}{syt}
  \DeclareMathOperator{\pyj}{pyj}
  \DeclareMathOperator{\dom}{dom}
  \DeclarePairedDelimiter{\abs}{\lvert}{\rvert}
  \DeclarePairedDelimiter{\norm}{\lVert}{\rVert}
  \DeclarePairedDelimiterXPP{\maxnorm}[1]{}{\lVert}{\rVert}{_\infty}{#1}
  \newcommand{\set}[1]{\left\lbrace#1\right\rbrace}
  \newcommand{\maxset}[1]{\max\set{#1}}
  \newcommand{\minset}[1]{\min\set{#1}}
  \newcommand{\powerset}[1]{\mathcal{P}\!\left(#1\right)}


\title{Lukuteoria I\\Harjoitus 5}
\author{}
\date{}

\begin{document}

\exerciseheading{Lukuteoria I}{Harjoitus 5}

\begin{exercise}
  \leavevmode\par
  \begin{enumerate}[label=(\alph*), topsep=0.35em]
    \item Laske \(\varphi(23)\), \(\varphi(105)\), \(\varphi(125)\) ja
      \(\varphi(700)\).
    \item Määritä lukujen 2, 3 ja 7 kertaluvut modulo 31.
  \end{enumerate}
\end{exercise}

\begin{solution}
  \leavevmode\par
  \begin{enumerate}[label=(\alph*), topsep=0.35em]
    \item Koska \(23\) on alkuluku, \(\varphi(23) = 23 - 1 = 22\).
      Koska \(105 = 3 \cdot 5 \cdot 7\) ja tekijät ovat pareittain
      keskenään jaottomia, funktion \(\varphi\) multiplikatiivisuudesta seuraa
      \[
        \varphi(105) = \varphi(3)\varphi(5)\varphi(7)
        = 2 \cdot 4 \cdot 6 = 48.
      \]
      Koska \(125 = 5^3\) ja \(700 = 2^2 \cdot 5^2 \cdot 7\), saadaan
      tehtävän 3 kaavalla
      \begin{align*}
        \varphi(125) &= 5^2(5 - 1) = 100, \\
        \varphi(700) &= 2(2 - 1) \cdot 5(5 - 1) \cdot (7 - 1) = 240.
      \end{align*}
    \item Lasketaan modulo \(31\). Olkoon \(d=\ord_{31}(a)\).
      Aloitetaan luvusta \(D=30=2\cdot 3\cdot 5\), joka on \(d\):n monikerta.
      Luvusta \(D\) voi poistaa alkutekijän \(q\), jos
      \(a^{D/q}\equiv 1\), eli jos \(D/q\) on edelleen \(d\):n monikerta.
      Kun \(q\) poistetaan, jatketaan luvulla \(D/q\). Saman alkutekijän
      voi poistaa uudelleen, jos se jakaa uutta lukua ja ehto yhä pätee.

      Jos \(D>d\), luvulla \(D/d>1\) on alkutekijä \(q\), jonka voi
      poistaa, sillä \(d\mid D/q\). Menettely päättyy siis, kun \(D=d\).
      \begin{itemize}
        \item \(2^5=32\equiv 1\). Koska \(5\) on alkuluku ja
          \(2\not\equiv 1\), saadaan \(\ord_{31}(2)=5\).
        \item Koska \(3^5=243\equiv -5\), eksponenteilla \(30/2\),
          \(30/3\) ja \(30/5\) saadaan
          \[
            3^{15}\equiv(-5)^3\equiv -1, \qquad
            3^{10}\equiv(-5)^2=25, \qquad
            3^6\equiv 3\cdot(-5)\equiv 16.
          \]
          Mikään näistä ei anna jäännöstä \(1\), joten \(\ord_{31}(3)=30\).
        \item Koska \(7^3=343\equiv 2\), saadaan \(7^{15}\equiv 2^5\equiv 1\).
          Eksponenteilla \(15/5\) ja \(15/3\) saadaan \(7^3\equiv 2\)
          ja \(7^5\equiv 2\cdot 18\equiv 5\), sillä \(7^2\equiv 18\).
          Siis \(\ord_{31}(7)=15\).
      \end{itemize}
  \end{enumerate}
\end{solution}

\begin{exercise}
  Olkoot \(m\), \(n\) ja \(k\) kokonaislukuja. Osoita, että
  \((k,mn) = 1\) jos ja vain jos \((k,m) = 1\) ja
  \((k,n) = 1\).
\end{exercise}

\begin{solution}
  \(\Rightarrow\) Oletetaan, että \((k,mn)=1\). Jokainen lukujen
  \(k\) ja \(m\) yhteinen jakaja on myös lukujen \(k\) ja \(mn\) yhteinen
  jakaja. Siis \((k,m)=1\). Samoin \((k,n)=1\).

  \(\Leftarrow\) Oletetaan, että \((k,m)=(k,n)=1\). Jos
  \((k,mn)\ne 1\), luvuilla \(k\) ja \(mn\) on yhteinen
  alkulukujakaja \(p\). Eukleideen lemman nojalla \(p\) jakaa joko luvun
  \(m\) tai luvun \(n\). Koska \(p\mid k\), tämä on ristiriidassa
  oletuksen \((k,m)=(k,n)=1\) kanssa. Siis \((k,mn)=1\).
\end{solution}

\begin{exercise}
  Olkoon kokonaisluvun \(n \ge 2\) alkutekijähajotelma
  \(n = p_1^{\alpha_1} p_2^{\alpha_2} \cdots p_k^{\alpha_k}\).
  Todista, että
  \[
    \varphi(n) = \prod_{j=1}^{k} p_j^{\alpha_j-1}(p_j-1).
  \]
\end{exercise}

\begin{solution}
  Olkoon \(p\) alkuluku ja \(\alpha\ge 1\) kokonaisluku. Jos
  \((a,p^\alpha)>1\), luvuilla \(a\) ja \(p^\alpha\) on yhteinen
  alkulukujakaja. Luvun \(p^\alpha\) ainoa alkutekijä on \(p\), joten
  \(p\mid a\). Kääntäen, jos \(p\mid a\), niin \(p\) jakaa molemmat
  luvut ja siten \((a,p^\alpha)>1\).

  Välillä \([1,p^\alpha]\) luvun \(p\) monikerrat ovat täsmälleen
  luvut \(pj\), missä \(j=1,\dots,p^{\alpha-1}\). Niitä on siis
  \(p^{\alpha-1}\), joten
  \[
    \varphi(p^\alpha) = p^\alpha - p^{\alpha-1} = p^{\alpha-1}(p-1).
  \]

  Koska \(p_1,\dots,p_k\) ovat eri alkulukuja, niiden potenssit
  \(p_1^{\alpha_1},\dots,p_k^{\alpha_k}\) ovat pareittain keskenään
  jaottomia. Soveltamalla \(\varphi\):n multiplikatiivisuutta toistuvasti
  saadaan
  \[
    \varphi(n) = \prod_{j=1}^k \varphi(p_j^{\alpha_j})
    = \prod_{j=1}^k p_j^{\alpha_j-1}(p_j-1).
  \]
\end{solution}

\begin{exercise}
  \leavevmode\par
  \begin{enumerate}[label=(\alph*), topsep=0.35em]
    \item Määritä kaikki primitiiviset juuret \(g\) modulo 19, kun
      \(1 \le g \le 18\).
    \item Määritä kaikki kokonaisluvut \(a\), joille \(1 \le a \le 30\) ja
      \(\operatorname{ord}_{31}(a) = 5\).
  \end{enumerate}
\end{exercise}

\begin{solution}
  \leavevmode\par
  \begin{enumerate}[label=(\alph*), topsep=0.35em]
    \item Luvun \(2\) kertaluku modulo \(19\) jakaa luvun \(18=2\cdot 3^2\).
      Tehtävän 1(b) tapaan tarkastellaan eksponentteja \(18/2\) ja \(18/3\).
      Koska \(2^9\equiv -1\) ja \(2^6\equiv 7\pmod{19}\), luvun \(2\)
      kertaluku on \(18\).

      Lauseen 71 mukaan potenssit \(2^0,\dots,2^{17}\)
      kattavat kaikki nollasta poikkeavat jäännösluokat modulo \(19\).
      Lemman 79 mukaan \(\ord_{19}(2^j)=18/(j,18)\), joten \(2^j\) on
      primitiivinen juuri, jos ja vain jos \((j,18)=1\).
      Sopivat eksponentit ovat \(1,5,7,11,13,17\), ja
      \[
        (2^1,2^5,2^7,2^{11},2^{13},2^{17})
        \equiv (2,13,14,15,3,10)\pmod{19}.
      \]
      Primitiiviset juuret ovat siis \(2,3,10,13,14,15\).

    \item Tehtävän 1(b) mukaan \(\ord_{31}(2)=5\). Lemman 79 mukaan
      \(\ord_{31}(2^j)=5/(j,5)\). Kun \(j=1,2,3,4\), saadaan siis
      luvut \(2,4,8,16\), joiden kertaluku on \(5\).

      Polynomilla \(X^5-1\in\mathbb{F}_{31}[X]\) on asteen nojalla enintään
      viisi juurta, ja \(1,2,4,8,16\) ovat sen juuria. Juurten kertaluvut
      jakavat alkuluvun \(5\), joten ratkaisut ovat täsmälleen juuret
      \(a\ne 1\).
  \end{enumerate}
\end{solution}

\begin{exercise}
  Olkoon \(g\) primitiivinen juuri modulo \(p\), missä \(p\) on alkuluku.
  Osoita, että \(g^j \equiv g^i \pmod{p}\) jos ja vain jos
  \(j \equiv i \pmod{p-1}\).
\end{exercise}

\begin{exercise}
  Olkoot \(p\) ja \(q\) parittomia alkulukuja, \(p = 2q + 1\) sekä \(a\)
  kokonaisluku, jolle \(1 < a < p - 1\). Osoita, että \(p - a^2\) on
  primitiivinen juuri modulo \(p\).
\end{exercise}

\end{document}
